\begin{problem}[Local geometry of a node]
Let $A=k[x,y]/(xy)$.  Compare the local rings at
\[
\mathfrak m=(x,y),\qquad
\mathfrak m_a=(x-a,y)\ (a\neq0),\qquad
\mathfrak p=(x).
\]
Determine their residue fields and explain which components are visible locally.
\end{problem}

\paragraph{Why this problem matters.}
This is the canonical example showing how localization separates components depending on the point.

\paragraph{Useful ideas before starting.}
At $\mathfrak m_a$, $x$ is a unit.  At $(x)$, $y$ is a unit.

\paragraph{Guided hints.}
Use $xy=0$ after inverting the appropriate coordinate.

\begin{solution}
At $\mathfrak m$, neither $x$ nor $y$ is inverted, so both remain nonzero zero divisors; $A_{\mathfrak m}$ sees both branches and has residue field $k$.  At $\mathfrak m_a$, $x$ is a unit, hence $y=0$, giving $A_{\mathfrak m_a}\cong k[x]_{(x-a)}$ with residue field $k$.  At $\mathfrak p=(x)$, $y$ is a unit and forces $x=0$, so $A_{\mathfrak p}\cong k(y)$, a field equal to its residue field.  It sees only the generic point of the $y$-axis.
\end{solution}

\paragraph{What happened mathematically.}
The local ring records exactly the components passing through the chosen point and forgets components removed by inversion.

\paragraph{Extensions.}
Repeat with three coordinate axes $k[x,y,z]/(xy,xz,yz)$.

\paragraph{Provenance.}
Expands the crossing-components local-ring examples and practice computation from legacy \emph{Theory of Commutative Algebra 12--13}.
